% GENERATED FILE. Edit canonical lessons, then run npm run generate:books.
% canonical-source-hash: fnv1a64:f6d68190804311fe
% canonical-lessons: ES-C530-arriesgado, ES-C530-riesgo, ES-C530-amenaza, ES-C530-ataque, ES-C530-atacar

\chapter{En peligro}
\label{ch:es-en-peligro}
\hlchaptermodality{530}

\begin{chapteropening}
\textbf{By the end of this chapter:} \emph{I can say that there was a risk or a threat, that something was risky, and that someone attacked.}

Complete the last lesson of chapter 530: I can say that there was a risk or a threat, that something was risky, and that someone attacked.
\end{chapteropening}

\section[arriesgado]{arriesgado — risky, daring}

\label{lesson:ES-C530-arriesgado}



\begin{quote}
Before the new one: say the Spanish for the rubble, then the Spanish for to destroy.
\end{quote}

\subsection*{You'll want to know: arriesgado}
\textbf{arriesgado} — \textquotedblleft{}risky, daring\textquotedblright{}.

\begin{quote}
Fue una idea muy arriesgada.
\end{quote}

\begin{quote}
\emph{It was a very risky idea.}
\end{quote}

\emph{Arriesgarse} is to take a risk.

\begin{cousinweb}
\emph{arriesgado} is built on the participle of \emph{arriesgar} 'to risk', from \emph{riesgo} 'a risk' (EtymDB). EtymDB takes \emph{riesgo} back to Italian \emph{rischio}; where that word came from is not settled.
\end{cousinweb}

\begin{grammarlens}[title={What you've built}]
One more word for telling, explaining and describing at B1.
\end{grammarlens}

\subsection*{Your turn}
\begin{itemize}
\raggedright
  \item \emph{Say it:} \emph{arriesgado}
  \item \emph{Say it:} \emph{arriesgado}, once more
  \item \emph{Say it:} a sentence of your own with \emph{arriesgado}
  \item \emph{Recall:} say the Spanish for the rubble, then the Spanish for to destroy, then say \emph{arriesgado} again
\end{itemize}

\begin{tcolorbox}[breakable,colback=teal!4,colframe=teal!35!black,title={Before you move on}]
What does arriesgado mean? (\textquotedblleft{}Risky, daring\textquotedblright{}.) Say it once more.
\end{tcolorbox}
\section[el riesgo]{el riesgo — the risk}

\label{lesson:ES-C530-riesgo}



\begin{quote}
Before the new one: say the Spanish for to destroy, then the Spanish for risky.
\end{quote}

\subsection*{You'll want to know: el riesgo}
\textbf{el riesgo} — \textquotedblleft{}the risk\textquotedblright{}.

\begin{quote}
Hay riesgo de incendio.
\end{quote}

\begin{quote}
\emph{There is a risk of fire.}
\end{quote}

\emph{Correr un riesgo} — to take a risk.

\begin{cousinweb}
The origin of \emph{riesgo} is not settled. EtymDB links it to old Spanish \emph{riesco} and Italian \emph{rischio}, but the older source of that word is still disputed.
\end{cousinweb}

\begin{grammarlens}[title={What you've built}]
One more word for telling, explaining and describing at B1.
\end{grammarlens}

\subsection*{Your turn}
\begin{itemize}
\raggedright
  \item \emph{Say it:} \emph{el riesgo}
  \item \emph{Say it:} \emph{el riesgo}, once more
  \item \emph{Say it:} a sentence of your own with \emph{el riesgo}
  \item \emph{Recall:} say the Spanish for to destroy, then the Spanish for risky, then say \emph{el riesgo} again
\end{itemize}

\begin{tcolorbox}[breakable,colback=teal!4,colframe=teal!35!black,title={Before you move on}]
What does el riesgo mean? (\textquotedblleft{}The risk\textquotedblright{}.) Say it once more.
\end{tcolorbox}
\section[la amenaza]{la amenaza — the threat}

\label{lesson:ES-C530-amenaza}



\begin{quote}
Before the new one: say the Spanish for risky, then the Spanish for the risk.
\end{quote}

\subsection*{You'll want to know: la amenaza}
\textbf{la amenaza} — \textquotedblleft{}the threat\textquotedblright{}.

\begin{quote}
Recibió una amenaza por teléfono.
\end{quote}

\begin{quote}
\emph{He received a threat over the phone.}
\end{quote}

\emph{Amenazar} is to threaten.

\begin{cousinweb}
\emph{amenaza} comes from spoken Latin \emph{*mināciā} 'a threat' (EtymDB), built on \emph{mināx} 'threatening', from \emph{minārī} 'to threaten'.
\end{cousinweb}

\begin{grammarlens}[title={What you've built}]
One more word for telling, explaining and describing at B1.
\end{grammarlens}

\subsection*{Your turn}
\begin{itemize}
\raggedright
  \item \emph{Say it:} \emph{la amenaza}
  \item \emph{Say it:} \emph{la amenaza}, once more
  \item \emph{Say it:} a sentence of your own with \emph{la amenaza}
  \item \emph{Recall:} say the Spanish for risky, then the Spanish for the risk, then say \emph{la amenaza} again
\end{itemize}

\begin{tcolorbox}[breakable,colback=teal!4,colframe=teal!35!black,title={Before you move on}]
What does la amenaza mean? (\textquotedblleft{}The threat\textquotedblright{}.) Say it once more.
\end{tcolorbox}
\section[el ataque]{el ataque — the attack}

\label{lesson:ES-C530-ataque}



\begin{quote}
Before the new one: say the Spanish for the risk, then the Spanish for the threat.
\end{quote}

\subsection*{You'll want to know: el ataque}
\textbf{el ataque} — \textquotedblleft{}the attack\textquotedblright{}.

\begin{quote}
El ataque ocurrió de noche.
\end{quote}

\begin{quote}
\emph{The attack happened at night.}
\end{quote}

\emph{Un ataque de risa} — a fit of laughter.

\begin{cousinweb}
\emph{ataque} is built on \emph{atacar} (EtymDB), which was borrowed from Italian \emph{attaccare} 'to attach, to join (battle)'.
\end{cousinweb}

\begin{grammarlens}[title={What you've built}]
One more word for telling, explaining and describing at B1.
\end{grammarlens}

\subsection*{Your turn}
\begin{itemize}
\raggedright
  \item \emph{Say it:} \emph{el ataque}
  \item \emph{Say it:} \emph{el ataque}, once more
  \item \emph{Say it:} a sentence of your own with \emph{el ataque}
  \item \emph{Recall:} say the Spanish for the risk, then the Spanish for the threat, then say \emph{el ataque} again
\end{itemize}

\begin{tcolorbox}[breakable,colback=teal!4,colframe=teal!35!black,title={Before you move on}]
What does el ataque mean? (\textquotedblleft{}The attack\textquotedblright{}.) Say it once more.
\end{tcolorbox}
\section[atacar]{atacar — to attack}

\label{lesson:ES-C530-atacar}



\begin{quote}
Before the new one: say the Spanish for the threat, then the Spanish for the attack.
\end{quote}

\subsection*{You'll want to know: atacar}
\textbf{atacar} — \textquotedblleft{}to attack\textquotedblright{}.

\begin{quote}
Un perro atacó al cartero.
\end{quote}

\begin{quote}
\emph{A dog attacked the postman.}
\end{quote}

It takes \emph{qu} before \emph{e}: \emph{ataqué}.

\begin{cousinweb}
\emph{atacar} was borrowed from Italian \emph{attaccare} 'to attach, to fasten' (EtymDB); to join battle was to fasten on to the enemy.
\end{cousinweb}

\begin{grammarlens}[title={What you've built}]
One more word for telling, explaining and describing at B1.
\end{grammarlens}

\subsection*{Your turn}
\begin{itemize}
\raggedright
  \item \emph{Say it:} \emph{atacar}
  \item \emph{Say it:} \emph{atacar}, once more
  \item \emph{Say it:} a sentence of your own with \emph{atacar}
  \item \emph{Recall:} say the Spanish for the threat, then the Spanish for the attack, then say \emph{atacar} again
\end{itemize}

\begin{tcolorbox}[breakable,colback=teal!4,colframe=teal!35!black,title={Before you move on}]
What does atacar mean? (\textquotedblleft{}To attack\textquotedblright{}.) Say it once more.
\end{tcolorbox}
